\documentclass[11pt]{article}
\usepackage[margin=1in]{geometry}
\usepackage{booktabs}
\title{Project 2: SCS-CN Runoff Model}
\author{Huang Qiwei}
\date{2026}
\begin{document}
\maketitle
\section*{Objective}
This project implements the SCS Curve Number method for estimating direct runoff from precipitation.
\section*{Method}
The model computes retention \(S = 25400/CN - 254\), initial abstraction \(I_a = 0.2S\), and runoff \(Q = (P-I_a)^2/(P-I_a+S)\) when rainfall exceeds initial abstraction.
\section*{Boundary Conditions}
\begin{tabular}{ll}
\toprule
Condition & Behavior \\
\midrule
\(P \le I_a\) & \(Q = 0\) \\
\(CN = 0\) & complete infiltration \\
\(CN = 100\) & \(Q = P\) \\
\(Q > P\) & clipped to \(P\) \\
\bottomrule
\end{tabular}
\section*{Outputs}
The script \texttt{sensitivity\_analysis.py} generates \texttt{scs\_cn\_sensitivity.png}, showing runoff sensitivity to curve number and precipitation depth.
\end{document}
